\chapter{Geometry basics}\label{ch:geom}
\index{geometry}\index{cross product}\index{convex hull}\index{Pick's theorem}\index{orientation test}\index{shoelace formula}

Contest geometry earns its reputation for fragility almost entirely
through one avoidable choice: writing \texttt{double} where an
integer would do. This chapter is the disciplined alternative.
After it you can run every standard test --- orientation, segment
intersection, point-in-polygon, convex hull, polygon area, lattice
points, closest pair --- on exact integer arithmetic, and you will
know which two or three degenerate cases actually decide whether
your code gets accepted.

\emph{This chapter is optional for the core route; it is the
foundation for USACO-G geometry and a handful of CF
1600--2100 problems.}

\section{The integer-only rule}\label{sec:intonly}
\index{integer geometry}\index{floating point!avoiding}
Input coordinates are integers, and nearly every geometric
predicate you need is a polynomial in them: orientation is a
quadratic form, containment a parity count, area a sum of
products. Doubles enter only if the \emph{answer} is real-valued
--- and even then, ``squared distance'' usually compares just as
well. The rule: \textbf{never compare two doubles for equality;
prefer an exact integer test and reach for \texttt{double} only
when printing a real number}. An \texttt{eps} band fixes nothing
--- it merely moves the failure to coordinates one order of
magnitude larger.

\begin{principle}\label{prin:geom}
Every primitive in \texttt{code/geom.h} returns an exact integer
or a fraction pair. If a comparison feels like it wants
\texttt{abs(a-b)<eps}, it almost always wants a cross product
instead.
\end{principle}

\section{Points, vectors, orientation}\label{sec:orient}
\index{orientation test}\index{cross product}
The workhorse is the 2D \emph{cross product}: for vectors
$\vec u=(u_x,u_y)$, $\vec v=(v_x,v_y)$,
\[
\vec u\times\vec v \;=\; u_x v_y-u_y v_x,
\]
the signed area of the parallelogram they span
(Figure~\ref{fig:orient}). Its sign is the \emph{orientation
test}: anchored at a common tail $o$,
$\texttt{orient}(o,a,b)=\texttt{cross}(a-o,\,b-o)$ is positive
when $b$ lies to the left of ray $o\to a$, negative to the right,
zero collinear. One determinant answers ``left, right, or on the
line'' --- the three-way predicate every later primitive builds
on.
\begin{figure}[ht]\centering
\begin{tikzpicture}[scale=1.05,font=\footnotesize]
\draw[-{stealth},slate] (0,0) -- (4.6,0) node[right]{$x$};
\draw[-{stealth},slate] (0,0) -- (0,3.2) node[above]{$y$};
\coordinate (O) at (0,0); \coordinate (A) at (3.6,1.2); \coordinate (B) at (1.4,2.6);
\draw[thick,ink,-{stealth}] (O) -- (A) node[below right]{$\vec u$};
\draw[thick,ink,-{stealth}] (O) -- (B) node[above left]{$\vec v$};
\draw[fill=amberbg,draw=amber,dashed] (O) -- (A) -- (3.6+1.4,1.2+2.6) -- (B) -- cycle;
\node[font=\footnotesize] at (2.6,1.7) {$\vec u \times \vec v$};
\node[align=left,anchor=west] at (5.0,2.2)
 {$>0$: $\vec v$ left of $\vec u$\\$<0$: right\\$=0$: collinear};
\end{tikzpicture}

\caption{The cross product as signed parallelogram area. Notice
the sign is the orientation test itself: swapping operands flips
it, which is why anchor order matters.}
\label{fig:orient}
\end{figure}
\lstinputlisting[linerange={11-25}]{code/geom.h}

\begin{example}[CSES \emph{Point Location Test}]\label{ex:loc}
For $t$ queries of three points $p,q,r$, print whether $r$ is
LEFT of, RIGHT of, or TOUCHing line $p\to q$.
\par\smallskip
\textbf{Strategy.} One orientation call per query; the sign of
$\texttt{orient}(p,q,r)$ \emph{is} the answer.
\par\smallskip
\begin{solution}
\lstinputlisting{code/point_location.cpp}
Real output (CSES sample):
\begin{lstlisting}[numbers=none]
LEFT
RIGHT
TOUCH
\end{lstlisting}
\end{solution}
\end{example}
\begin{exercise}[easy]\label{exr:signflip}
Show $\texttt{orient}(o,a,b)=-\texttt{orient}(o,b,a)$ from the
determinant. Which order does your library use, and where would
a flipped sign hide in a hull?
\end{exercise}

\section{Segment intersection}\label{sec:segint}
\index{segment intersection}
Two segments $ab$, $cd$ intersect exactly when $c$ and $d$ lie on
\emph{opposite sides} of line $ab$ \emph{and} $a,b$ lie on
opposite sides of line $cd$ --- two orientation tests each way.
But four touching cases matter just as much: endpoint touching,
collinear overlap, an endpoint \emph{on} the other segment, and
the collinear-but-disjoint miss. Handle every zero-orientation
case \emph{before} the strict sign test
(Figure~\ref{fig:segint}).
\begin{figure}[ht]\centering
% fig79-segint.tex -- Figure 15.x: the four segment-intersection cases.
% fstyle.tex -- shared TikZ styles for the figures of this book (palette from lifeos.sty).
\tikzset{%
  cell/.style={draw=ink, minimum width=7mm, minimum height=6.5mm, inner sep=0pt, font=\footnotesize, fill=white},
  cellh/.style={cell, fill=accentsoft, draw=accent},
  cella/.style={cell, fill=amberbg, draw=amber},
  cellg/.style={cell, fill=goodbg, draw=good},
  cellx/.style={cell, fill=graybg, draw=slate, text=slate},
  vtx/.style={circle, draw=ink, minimum size=6.5mm, inner sep=0pt, font=\footnotesize, fill=white},
  vtxh/.style={vtx, draw=accent, fill=accentsoft},
  vtxa/.style={vtx, draw=amber, fill=amberbg},
  vtxg/.style={vtx, draw=good, fill=goodbg},
  vtxd/.style={vtx, draw=slate, fill=graybg, text=slate},
  lab/.style={font=\footnotesize, text=slate},
  labA/.style={font=\footnotesize, text=accent},
  labM/.style={font=\footnotesize, text=amber},
  idx/.style={font=\tiny, text=slate},
  arr/.style={-{stealth}, thick, ink},
  arrA/.style={-{stealth}, thick, accent},
  arrM/.style={-{stealth}, thick, amber},
  arrG/.style={-{stealth}, thick, good},
  edge/.style={thick, ink},
  edgeA/.style={thick, accent},
  edgeM/.style={thick, amber},
}%
% \cellrow{prefix}{x0}{y}{style}{v0,v1,...}: a row of cells named prefix0, prefix1, ... (pitch 0.75cm)
\providecommand{\cellrow}[5]{\foreach \v [count=\i from 0] in {#5} {\node[#4] (#1\i) at (#2+\i*0.75,#3) {\v};}}%
% \cellidx{prefix}{n}{start}: index labels under cells prefix0 .. prefix(n-1)
\providecommand{\cellidx}[3]{\foreach \i in {0,...,\numexpr#2-1\relax} {\node[idx, below=0pt of #1\i] {\the\numexpr\i+#3\relax};}}%
% \fitfig{<tikzpicture>}: shrink the picture to \linewidth only when it would overflow.
\ifdefined\fitbox\else\newsavebox{\fitbox}\fi%
\providecommand{\fitfig}[1]{\sbox{\fitbox}{#1}\ifdim\wd\fitbox>\linewidth\resizebox{\linewidth}{!}{\usebox{\fitbox}}\else\usebox{\fitbox}\fi}%

\begin{tikzpicture}[font=\footnotesize]
\tikzset{sa/.style={ultra thick, accent}, sb/.style={ultra thick, amber}}
% proper
\begin{scope}
  \draw[sa] (0,0)--(2.6,1.8); \draw[sb] (0.3,1.8)--(2.3,0.2);
  \fill[ink] (1.32,0.92) circle (2pt);
  \node[lab] at (1.3,-0.7) {proper crossing};
\end{scope}
% touch
\begin{scope}[xshift=3.8cm]
  \draw[sa] (0,0)--(2.6,1.1); \draw[sb] (2.6,1.1)--(1.6,2.2);
  \fill[ink] (2.6,1.1) circle (2pt);
  \node[lab] at (1.3,-0.7) {endpoint touches};
\end{scope}
% overlap
\begin{scope}[xshift=7.6cm]
  \draw[sa] (0,0.2)--(2.4,1.1); \draw[sb] ($(0,0.2)!0.45!(2.4,1.1)$)--($(0,0.2)!1.25!(2.4,1.1)$);
  \node[lab] at (1.3,-0.7) {collinear overlap};
\end{scope}
% miss
\begin{scope}[xshift=11.4cm]
  \draw[sa] (0,0)--(2.2,1.5); \draw[sb] (0.5,2.2)--(2.4,0.7);
  \node[lab] at (1.3,-0.7) {no intersection};
\end{scope}
\end{tikzpicture}

\caption{The four verdicts. Notice ``touch'' and ``overlap'' are
both YES for CSES --- the sign test alone sees neither, because
a zero kills the strict comparison.}
\label{fig:segint}
\end{figure}
\lstinputlisting[linerange={27-38}]{code/geom.h}
\begin{example}[CSES \emph{Line Segment Intersection}]\label{ex:seg}
$t$ queries of two segments; print YES iff they intersect
(touching counts).
\par\smallskip
\textbf{Strategy.} Four orientation calls; test the
endpoint-on-segment cases first, bail on any remaining zero,
then compare signs both ways.
\par\smallskip
\begin{solution}
\lstinputlisting{code/line_seg_int.cpp}
Real output (CSES sample):
\begin{lstlisting}[numbers=none]
NO
YES
YES
YES
YES
\end{lstlisting}
\end{solution}
\end{example}
\begin{exercise}[easy]\label{exr:cases}
The sample contains an endpoint touch and a collinear overlap.
Identify which lines they are, and argue why the strict sign
test alone returns NO for each.
\end{exercise}

\section{Convex hull}\label{sec:hull}
\index{convex hull}\index{monotone chain}
The \emph{convex hull} of a point set is the smallest convex
polygon containing it --- the rubber band around pegs. Andrew's
\emph{monotone chain} builds it in $O(n\log n)$ with no angle
sort: sort by $(x,y)$, sweep left-to-right for the lower chain
and right-to-left for the upper, popping the last vertex
whenever the new point makes a non-left turn
(Figure~\ref{fig:hull}).
\begin{figure}[ht]\centering
% fig80-hullbuild.tex -- Figure 15.x: monotone chain on {(2,1),(2,5),(3,3),(4,3),(4,4),(6,3)}; hull (2,1),(6,3),(2,5).
% fstyle.tex -- shared TikZ styles for the figures of this book (palette from lifeos.sty).
\tikzset{%
  cell/.style={draw=ink, minimum width=7mm, minimum height=6.5mm, inner sep=0pt, font=\footnotesize, fill=white},
  cellh/.style={cell, fill=accentsoft, draw=accent},
  cella/.style={cell, fill=amberbg, draw=amber},
  cellg/.style={cell, fill=goodbg, draw=good},
  cellx/.style={cell, fill=graybg, draw=slate, text=slate},
  vtx/.style={circle, draw=ink, minimum size=6.5mm, inner sep=0pt, font=\footnotesize, fill=white},
  vtxh/.style={vtx, draw=accent, fill=accentsoft},
  vtxa/.style={vtx, draw=amber, fill=amberbg},
  vtxg/.style={vtx, draw=good, fill=goodbg},
  vtxd/.style={vtx, draw=slate, fill=graybg, text=slate},
  lab/.style={font=\footnotesize, text=slate},
  labA/.style={font=\footnotesize, text=accent},
  labM/.style={font=\footnotesize, text=amber},
  idx/.style={font=\tiny, text=slate},
  arr/.style={-{stealth}, thick, ink},
  arrA/.style={-{stealth}, thick, accent},
  arrM/.style={-{stealth}, thick, amber},
  arrG/.style={-{stealth}, thick, good},
  edge/.style={thick, ink},
  edgeA/.style={thick, accent},
  edgeM/.style={thick, amber},
}%
% \cellrow{prefix}{x0}{y}{style}{v0,v1,...}: a row of cells named prefix0, prefix1, ... (pitch 0.75cm)
\providecommand{\cellrow}[5]{\foreach \v [count=\i from 0] in {#5} {\node[#4] (#1\i) at (#2+\i*0.75,#3) {\v};}}%
% \cellidx{prefix}{n}{start}: index labels under cells prefix0 .. prefix(n-1)
\providecommand{\cellidx}[3]{\foreach \i in {0,...,\numexpr#2-1\relax} {\node[idx, below=0pt of #1\i] {\the\numexpr\i+#3\relax};}}%
% \fitfig{<tikzpicture>}: shrink the picture to \linewidth only when it would overflow.
\ifdefined\fitbox\else\newsavebox{\fitbox}\fi%
\providecommand{\fitfig}[1]{\sbox{\fitbox}{#1}\ifdim\wd\fitbox>\linewidth\resizebox{\linewidth}{!}{\usebox{\fitbox}}\else\usebox{\fitbox}\fi}%

\begin{tikzpicture}[font=\footnotesize, x=1.0cm, y=0.85cm]
\draw[slate!30] (1,0) grid (7,6);
\foreach \x in {1,...,6} \node[idx] at (\x+0.5-0.5,-0.25) {\x}; \foreach \y in {1,...,5} \node[idx, anchor=east] at (0.8,\y) {\y};
\fill[accent!10] (2,1)--(6,3)--(2,5)--cycle;
\draw[accent, very thick] (2,1)--(6,3) node[midway, below right, labA] {lower chain};
\draw[good, very thick] (6,3)--(2,5) node[pos=0.78, above right=-1pt, text=good, font=\footnotesize] {upper chain};
\draw[slate, very thick] (2,5)--(2,1);
\foreach \x/\y/\sty in {2/1/ink, 2/5/ink, 6/3/ink} \fill[\sty] (\x,\y) circle (3pt);
\foreach \x/\y in {3/3, 4/3} \fill[slate] (\x,\y) circle (2.2pt);
\fill[amber] (4,4) circle (2.6pt);
\node[font=\scriptsize, anchor=south west] at (2,1) {(2,1)}; \node[font=\scriptsize, anchor=south west] at (2,5) {(2,5)}; \node[font=\scriptsize, anchor=south west] at (6,3) {(6,3)};
\node[font=\scriptsize, text=slate, anchor=north west] at (3,3) {(3,3)}; \node[font=\scriptsize, text=slate, anchor=north west] at (4,3) {(4,3)};
\node[font=\scriptsize, text=amber, anchor=north west] at (4.05,3.95) {(4,4) collinear};
\node[lab, anchor=west, align=left, text width=5.2cm] at (7.4,4.6) {sort by $(x,y)$; sweep left to right for the lower chain, right to left for the upper chain};
\node[lab, anchor=west, align=left, text width=5.2cm] at (7.4,2.8) {pop the last vertex whenever the new point makes a non-left turn ($\text{cross}\le 0$)};
\node[labM, anchor=west, align=left, text width=5.2cm] at (7.4,1.1) {$(4,4)$ lies exactly on the segment $(6,3)$--$(2,5)$, so it is popped; interior points never survive};
\end{tikzpicture}

\caption{Monotone chain on the CSES sample. Notice the dashed
chord: $(3,3)$ and $(4,3)$ never reach the hull --- every pop is
one orientation test.}
\label{fig:hull}
\end{figure}
The pop condition is the one design decision: pop on
\texttt{orient < 0} keeps collinear boundary points on the hull
(CSES wants them); pop on \texttt{orient <= 0} keeps only extreme
vertices. Read the problem for which hull it defines.
\lstinputlisting[linerange={76-96}]{code/geom.h}
\begin{example}[CSES \emph{Convex Hull}]\label{ex:hull}
Print the hull's vertex count, then its vertices (boundary
points included).
\par\smallskip
\textbf{Strategy.} Sort, two monotone passes with strict
\texttt{<0} pops so collinear boundary points survive. Any
cyclic order of the same vertices is accepted.
\par\smallskip
\begin{solution}
\lstinputlisting{code/convex_hull.cpp}
Real output (CSES sample):
\begin{lstlisting}[numbers=none]
4
2 1
6 3
4 4
2 5
\end{lstlisting}
\end{solution}
\end{example}

\section{Polygon measures: area, lattice points, containment}\label{sec:poly}
\index{shoelace formula}\index{Pick's theorem}\index{point in polygon}
The \emph{shoelace} sums cross products of consecutive vertices:
$\texttt{area2}=\sum_i x_i y_{i+1}-x_{i+1}y_i$, exactly twice the
signed area --- keep it doubled and integer. \emph{Pick's
theorem} ties area to the lattice:
$A = I + B/2 - 1$, where $I$ counts interior lattice points and
$B$ the boundary ones, and each edge $(d_x,d_y)$ contributes
$\gcd(|d_x|,|d_y|)$ to $B$ (Figure~\ref{fig:pick}).
\begin{figure}[ht]\centering
% fig81-pick.tex -- Figure 15.x: Pick's theorem on polygon (1,1)-(5,3)-(3,5)-(1,4): A=9, B=8, I=6 (A = I + B/2 - 1).
% fstyle.tex -- shared TikZ styles for the figures of this book (palette from lifeos.sty).
\tikzset{%
  cell/.style={draw=ink, minimum width=7mm, minimum height=6.5mm, inner sep=0pt, font=\footnotesize, fill=white},
  cellh/.style={cell, fill=accentsoft, draw=accent},
  cella/.style={cell, fill=amberbg, draw=amber},
  cellg/.style={cell, fill=goodbg, draw=good},
  cellx/.style={cell, fill=graybg, draw=slate, text=slate},
  vtx/.style={circle, draw=ink, minimum size=6.5mm, inner sep=0pt, font=\footnotesize, fill=white},
  vtxh/.style={vtx, draw=accent, fill=accentsoft},
  vtxa/.style={vtx, draw=amber, fill=amberbg},
  vtxg/.style={vtx, draw=good, fill=goodbg},
  vtxd/.style={vtx, draw=slate, fill=graybg, text=slate},
  lab/.style={font=\footnotesize, text=slate},
  labA/.style={font=\footnotesize, text=accent},
  labM/.style={font=\footnotesize, text=amber},
  idx/.style={font=\tiny, text=slate},
  arr/.style={-{stealth}, thick, ink},
  arrA/.style={-{stealth}, thick, accent},
  arrM/.style={-{stealth}, thick, amber},
  arrG/.style={-{stealth}, thick, good},
  edge/.style={thick, ink},
  edgeA/.style={thick, accent},
  edgeM/.style={thick, amber},
}%
% \cellrow{prefix}{x0}{y}{style}{v0,v1,...}: a row of cells named prefix0, prefix1, ... (pitch 0.75cm)
\providecommand{\cellrow}[5]{\foreach \v [count=\i from 0] in {#5} {\node[#4] (#1\i) at (#2+\i*0.75,#3) {\v};}}%
% \cellidx{prefix}{n}{start}: index labels under cells prefix0 .. prefix(n-1)
\providecommand{\cellidx}[3]{\foreach \i in {0,...,\numexpr#2-1\relax} {\node[idx, below=0pt of #1\i] {\the\numexpr\i+#3\relax};}}%
% \fitfig{<tikzpicture>}: shrink the picture to \linewidth only when it would overflow.
\ifdefined\fitbox\else\newsavebox{\fitbox}\fi%
\providecommand{\fitfig}[1]{\sbox{\fitbox}{#1}\ifdim\wd\fitbox>\linewidth\resizebox{\linewidth}{!}{\usebox{\fitbox}}\else\usebox{\fitbox}\fi}%

\begin{tikzpicture}[font=\footnotesize, x=0.9cm, y=0.9cm]
\draw[slate!35] (0,0) grid (6,6);
\fill[accent!10] (1,1)--(5,3)--(3,5)--(1,4)--cycle; \draw[ink, very thick] (1,1)--(5,3)--(3,5)--(1,4)--cycle;
\fill[accent] (2,2) circle (3pt);
\fill[accent] (2,3) circle (3pt);
\fill[accent] (2,4) circle (3pt);
\fill[accent] (3,3) circle (3pt);
\fill[accent] (3,4) circle (3pt);
\fill[accent] (4,3) circle (3pt);
\fill[good] (1,1) circle (2.6pt);
\fill[good] (1,2) circle (2.6pt);
\fill[good] (1,3) circle (2.6pt);
\fill[good] (1,4) circle (2.6pt);
\fill[good] (3,2) circle (2.6pt);
\fill[good] (3,5) circle (2.6pt);
\fill[good] (4,4) circle (2.6pt);
\fill[good] (5,3) circle (2.6pt);
\node[lab, anchor=west, align=left] at (6.8,5.0) {area $A=9$ (shoelace: $2A=18$)};
\node[font=\footnotesize, text=good, anchor=west] at (6.8,4.0) {$\bullet$ boundary points $B=8$};
\node[labA, anchor=west] at (6.8,3.2) {$\bullet$ interior points $I=6$};
\node[labM, anchor=west, align=left] at (6.8,2.0) {$I=A-\dfrac B2+1=9-4+1=6$ \checkmark};
\end{tikzpicture}

\caption{Pick's theorem on the CSES sample. Notice $B$ counts
edge \emph{segments'} interior lattice points too --- the gcd
adds every lattice step along an edge.}
\label{fig:pick}
\end{figure}
\begin{example}[CSES \emph{Polygon Area}]\label{ex:area}
Print twice the area of a simple polygon.
\par\smallskip
\textbf{Strategy.} Shoelace, take the absolute value.
\par\smallskip
\begin{solution}
\lstinputlisting{code/poly_area.cpp}
Real output (CSES sample):
\begin{lstlisting}[numbers=none]
16
\end{lstlisting}
\end{solution}
\end{example}
\begin{example}[CSES \emph{Polygon Lattice Points}]\label{ex:pick}
Print interior then boundary lattice-point counts.
\par\smallskip
\textbf{Strategy.} $B$ from the gcd sum; invert Pick:
$I=(\texttt{area2}-B+2)/2$.
\par\smallskip
\begin{solution}
\lstinputlisting{code/poly_lattice.cpp}
Real output (CSES sample):
\begin{lstlisting}[numbers=none]
6 8
\end{lstlisting}
\end{solution}
\end{example}
Containment is the \emph{ray-casting parity} test: shoot a
horizontal ray, count polygon edges it crosses --- odd means
inside. The half-open rule (count an edge only when exactly one
endpoint is strictly above the ray) is what stops vertex touches
from double-counting (Figure~\ref{fig:raycast}); boundary points
are caught by an explicit on-segment check first.
\begin{figure}[ht]\centering
% fig82-raycast.tex -- Figure 15.x: ray casting. Polygon (0,0),(4,0),(5,3),(2,4),(0,3); ray to the right at y=2.
% fstyle.tex -- shared TikZ styles for the figures of this book (palette from lifeos.sty).
\tikzset{%
  cell/.style={draw=ink, minimum width=7mm, minimum height=6.5mm, inner sep=0pt, font=\footnotesize, fill=white},
  cellh/.style={cell, fill=accentsoft, draw=accent},
  cella/.style={cell, fill=amberbg, draw=amber},
  cellg/.style={cell, fill=goodbg, draw=good},
  cellx/.style={cell, fill=graybg, draw=slate, text=slate},
  vtx/.style={circle, draw=ink, minimum size=6.5mm, inner sep=0pt, font=\footnotesize, fill=white},
  vtxh/.style={vtx, draw=accent, fill=accentsoft},
  vtxa/.style={vtx, draw=amber, fill=amberbg},
  vtxg/.style={vtx, draw=good, fill=goodbg},
  vtxd/.style={vtx, draw=slate, fill=graybg, text=slate},
  lab/.style={font=\footnotesize, text=slate},
  labA/.style={font=\footnotesize, text=accent},
  labM/.style={font=\footnotesize, text=amber},
  idx/.style={font=\tiny, text=slate},
  arr/.style={-{stealth}, thick, ink},
  arrA/.style={-{stealth}, thick, accent},
  arrM/.style={-{stealth}, thick, amber},
  arrG/.style={-{stealth}, thick, good},
  edge/.style={thick, ink},
  edgeA/.style={thick, accent},
  edgeM/.style={thick, amber},
}%
% \cellrow{prefix}{x0}{y}{style}{v0,v1,...}: a row of cells named prefix0, prefix1, ... (pitch 0.75cm)
\providecommand{\cellrow}[5]{\foreach \v [count=\i from 0] in {#5} {\node[#4] (#1\i) at (#2+\i*0.75,#3) {\v};}}%
% \cellidx{prefix}{n}{start}: index labels under cells prefix0 .. prefix(n-1)
\providecommand{\cellidx}[3]{\foreach \i in {0,...,\numexpr#2-1\relax} {\node[idx, below=0pt of #1\i] {\the\numexpr\i+#3\relax};}}%
% \fitfig{<tikzpicture>}: shrink the picture to \linewidth only when it would overflow.
\ifdefined\fitbox\else\newsavebox{\fitbox}\fi%
\providecommand{\fitfig}[1]{\sbox{\fitbox}{#1}\ifdim\wd\fitbox>\linewidth\resizebox{\linewidth}{!}{\usebox{\fitbox}}\else\usebox{\fitbox}\fi}%

\begin{tikzpicture}[font=\footnotesize, x=1.1cm, y=0.95cm]
\draw[ink, very thick, fill=accent!8] (0,0)--(4,0)--(5,3)--(2,4)--(0,3)--cycle;
% inside point q
\fill[accent] (1.5,2) circle (2.6pt) node[below left, labA] {$q$};
\draw[arrA, dashed] (1.5,2)--(6.3,2);
\fill[amber] (4.667,2) circle (2.6pt) node[below right=1pt, labM] {};
% outside point p
\fill[good] (5.4,1.0) circle (2.6pt) node[below, text=good, font=\footnotesize] {$p$};
\draw[arrG, dashed] (5.4,1.0)--(6.3,1.0);
\node[font=\footnotesize, text=good, anchor=west] at (6.4,1.0) {0 crossings: outside};
\node[font=\footnotesize, text=accent, anchor=west] at (6.4,2.0) {1 crossing (odd): inside};
\node[lab, anchor=west, align=left, text width=5.6cm] at (6.4,3.4) {cast a ray from the point and count edges it crosses};
\node[labM, anchor=west, align=left, text width=5.6cm] at (6.4,-0.2) {a vertex on the ray is counted once via the half-open rule: edge $(a,b)$ counts if $a_y\le y<b_y$ or $b_y\le y<a_y$};
\end{tikzpicture}

\caption{Parity ray cast. Notice which edges count: the rule is
``one endpoint strictly above the ray,'' never both or neither.}
\label{fig:raycast}
\end{figure}
\lstinputlisting[linerange={62-73}]{code/geom.h}
\begin{example}[CSES \emph{Point in Polygon}]\label{ex:pip}
$m$ queries against an $n$-gon: INSIDE, OUTSIDE, or BOUNDARY.
\par\smallskip
\textbf{Strategy.} On-segment check first, then the parity flip
per crossing edge.
\par\smallskip
\begin{solution}
\lstinputlisting{code/point_in_poly.cpp}
Real output (CSES sample):
\begin{lstlisting}[numbers=none]
INSIDE
OUTSIDE
BOUNDARY
\end{lstlisting}
\end{solution}
\end{example}
\begin{exercise}[medium]\label{exr:pickinv}
A lattice polygon has $I=6$, $B=8$. Find $A$ --- then exhibit a
\emph{different} polygon with the same $(I,B)$ and convince
yourself Pick still applies.
\end{exercise}

\section{Distances}\label{sec:dist}
\index{distance!squared}
Distances stay squared. $\texttt{dist2}(a,b)$ is exact, ordering
preserved, and the closest-pair problem wants it anyway.
Minimum pairwise distance among $n$ points: sort by $x$, divide
at the middle, recurse, then examine only the strip of width
$d=\min(d_L,d_R)$ around the cut --- scanning the strip in
$y$-order touches a constant number of candidates per point, for
$O(n\log n)$ overall (Figure~\ref{fig:strip}).
\begin{figure}[ht]\centering
% fig83-strip.tex -- Figure 15.x: closest-pair strip around the dividing line x = m.
% fstyle.tex -- shared TikZ styles for the figures of this book (palette from lifeos.sty).
\tikzset{%
  cell/.style={draw=ink, minimum width=7mm, minimum height=6.5mm, inner sep=0pt, font=\footnotesize, fill=white},
  cellh/.style={cell, fill=accentsoft, draw=accent},
  cella/.style={cell, fill=amberbg, draw=amber},
  cellg/.style={cell, fill=goodbg, draw=good},
  cellx/.style={cell, fill=graybg, draw=slate, text=slate},
  vtx/.style={circle, draw=ink, minimum size=6.5mm, inner sep=0pt, font=\footnotesize, fill=white},
  vtxh/.style={vtx, draw=accent, fill=accentsoft},
  vtxa/.style={vtx, draw=amber, fill=amberbg},
  vtxg/.style={vtx, draw=good, fill=goodbg},
  vtxd/.style={vtx, draw=slate, fill=graybg, text=slate},
  lab/.style={font=\footnotesize, text=slate},
  labA/.style={font=\footnotesize, text=accent},
  labM/.style={font=\footnotesize, text=amber},
  idx/.style={font=\tiny, text=slate},
  arr/.style={-{stealth}, thick, ink},
  arrA/.style={-{stealth}, thick, accent},
  arrM/.style={-{stealth}, thick, amber},
  arrG/.style={-{stealth}, thick, good},
  edge/.style={thick, ink},
  edgeA/.style={thick, accent},
  edgeM/.style={thick, amber},
}%
% \cellrow{prefix}{x0}{y}{style}{v0,v1,...}: a row of cells named prefix0, prefix1, ... (pitch 0.75cm)
\providecommand{\cellrow}[5]{\foreach \v [count=\i from 0] in {#5} {\node[#4] (#1\i) at (#2+\i*0.75,#3) {\v};}}%
% \cellidx{prefix}{n}{start}: index labels under cells prefix0 .. prefix(n-1)
\providecommand{\cellidx}[3]{\foreach \i in {0,...,\numexpr#2-1\relax} {\node[idx, below=0pt of #1\i] {\the\numexpr\i+#3\relax};}}%
% \fitfig{<tikzpicture>}: shrink the picture to \linewidth only when it would overflow.
\ifdefined\fitbox\else\newsavebox{\fitbox}\fi%
\providecommand{\fitfig}[1]{\sbox{\fitbox}{#1}\ifdim\wd\fitbox>\linewidth\resizebox{\linewidth}{!}{\usebox{\fitbox}}\else\usebox{\fitbox}\fi}%

\begin{tikzpicture}[font=\footnotesize, x=1cm, y=0.9cm]
\fill[accent!10] (2.2,0) rectangle (4.2,4.6);
\draw[dashed, slate] (3.2,-0.2)--(3.2,4.8) node[above, lab] {$x=m$};
\draw[accent, thick] (2.2,0) rectangle (4.2,4.6);
\foreach \x/\y in {0.3/1.2,0.8/3.4,1.4/2.0,1.7/4.1,2.5/1.0,3.0/3.0,3.5/2.2,4.6/3.6,5.2/1.4,5.7/2.8} \fill[slate] (\x,\y) circle (2pt);
\foreach \x/\y in {2.5/3.5,2.9/2.6,3.5/3.9,3.8/1.9} \fill[accent] (\x,\y) circle (2.4pt);
\draw[amber, thick, dashed] (2.2,2.6) rectangle (4.2,3.71);
\node[font=\scriptsize, text=amber, anchor=west] at (4.25,3.15) {height $d$};
\draw[<->, amber, thick] (2.2,-0.5)--(3.2,-0.5) node[midway, below, labM] {$d$}; \draw[<->, amber, thick] (3.2,-0.5)--(4.2,-0.5) node[midway, below, labM] {$d$};
\node[lab, anchor=west, align=left, text width=5.6cm] at (6.5,4.3) {$d=\min(d_L,d_R)$ from the two halves};
\node[lab, anchor=west, align=left, text width=5.2cm] at (6.5,3.3) {only points within $d$ of the line (blue) can beat $d$};
\node[labA, anchor=west, align=left, text width=5.2cm] at (6.5,1.4) {scanned in $y$-order, a point only checks the few strip points within $d$ above it: past a gap of $d$ no later point can improve $d$};
\end{tikzpicture}

\caption{Closest pair's merge strip. Notice why $y$-order
suffices: past a gap of $d$ in $y$ no later point can improve
$d$.}
\label{fig:strip}
\end{figure}
\begin{example}[CSES \emph{Minimum Euclidean Distance}]\label{ex:mindist}
Print the smallest squared distance between two of $n$ points,
$n\le 2\cdot10^5$.
\par\smallskip
\textbf{Strategy.} Divide and conquer; never a square root ---
the output is $d^2$.
\par\smallskip
\begin{solution}
\lstinputlisting{code/min_dist.cpp}
Real output (CSES sample):
\begin{lstlisting}[numbers=none]
2
\end{lstlisting}
\end{solution}
\end{example}
\begin{exercise}[medium]\label{exr:strip}
Why is a constant-size inner scan correct? Bound the number of
strip points inside a $d\times d$ box when every pair in a half
is at least $d$ apart.
\end{exercise}

\section{Common mistakes}
\index{common mistakes}
\begin{itemize}
\item \textbf{Epsilon equality.} \texttt{abs(a-b)<eps} is not an
integer test; it fails silently one magnitude up. Use cross
products.
\item \textbf{Sign-first segment intersection.} The strict
double-orientation test misses every touching case --- endpoint
touch, collinear overlap --- which CSES counts as YES. Check
zero orientations first.
\item \textbf{Popping collinear vertices.} \texttt{orient<=0}
pops strip the boundary points CSES \emph{Convex Hull} asks
for; \texttt{<0} keeps them. Match the problem's definition.
\item \textbf{Vertex-touching ray casts.} Counting every edge
the ray touches double-counts polygon vertices on the ray; the
half-open $y$ rule fixes parity.
\item \textbf{Overflowing \texttt{int}.} Coordinates up to
$10^9$ make cross products reach $10^{18}$ --- \texttt{long
long} or nothing, and keep squared distance out of
\texttt{sqrt} until the answer truly needs decimals.
\end{itemize}

\section*{Summary}
\begin{tabularx}{\linewidth}{@{}lXl@{}}
\textbf{Need} & \textbf{Primitive} & \textbf{Cost}\\
\midrule
left/right/collinear & \texttt{orient(o,a,b)} sign & $O(1)$\\
segments meet & \texttt{seg\_int}: zeros first, then both sign tests & $O(1)$\\
twice signed area & shoelace \texttt{area2} & $O(n)$\\
inside/outside/boundary & on-seg check + ray parity & $O(n)$/$q$\\
lattice $I,B$ & gcd boundary sum, invert Pick & $O(n)$\\
hull (with boundary) & monotone chain, pop on $<0$ & $O(n\log n)$\\
closest pair & divide \& conquer + $y$-strip & $O(n\log n)$\\
\end{tabularx}

\section*{Exercises}
\begin{exercise}[easy] CSES \emph{Point Location Test} --- re-implement \texttt{orient} unaided. \end{exercise}
\begin{exercise}[easy] USACO-B \emph{Square Pasture} --- axis-aligned bounding square, pure integers. \end{exercise}
\begin{exercise}[easy] USACO-B \emph{White Sheet} --- rectangle overlap area. \end{exercise}
\begin{exercise}[easy] Kattis \emph{Segment Distance} --- your \texttt{seg\_dist2} fraction, printed as a real at the end. \end{exercise}
\begin{exercise}[easy] CF \emph{Beru-taxi} --- nearest point; compare squared distances only. \end{exercise}
\begin{exercise}[medium] CSES \emph{Line Segment Intersection} --- full edge-case drill. \end{exercise}
\begin{exercise}[medium] CSES \emph{Polygon Lattice Points} --- gcd sum + Pick. \end{exercise}
\begin{exercise}[medium] Kattis \emph{Convex Hull} --- monotone chain; compare against the collinear-keeping variant. \end{exercise}
\begin{exercise}[medium] Kattis \emph{Robot Protection} --- hull then shoelace area. \end{exercise}
\begin{exercise}[medium] CF \emph{Water Lily} --- one right triangle, one equation. \end{exercise}
\begin{exercise}[medium] CSES \emph{Point in Polygon} --- parity + boundary handling together. \end{exercise}
\begin{exercise}[hard] USACO-S \emph{Cow Steeplechase II} --- all-pairs segment intersection with a twist. \end{exercise}
\begin{exercise}[hard] USACO-P \emph{Cow Curling} --- hull maintenance inside a DP. \end{exercise}
\begin{exercise}[hard] Kattis \emph{Closest Pair} --- your strip, pushed to adversarial inputs. \end{exercise}
\begin{exercise}[hard] Derive the \emph{rotating calipers} diameter of a convex polygon from \texttt{orient}, then implement it. \end{exercise}

\chapternotes

\section*{Further reading}
\begin{itemize}
\item A.~Laaksonen, \emph{Competitive Programmer's Handbook},
ch.~29 ``Geometry'', \url{https://cses.fi/book/book.pdf} ---
free PDF; the cross-product-first presentation this chapter
tracks.
\item cp-algorithms, \emph{Basic Geometry} --- the same
integer-first discipline with more primitives.
\newline\url{https://cp-algorithms.com/geometry/basic-geometry.html}
\item cp-algorithms, \emph{Convex hull construction} ---
monotone chain plus the Graham-scan alternative.
\newline\url{https://cp-algorithms.com/geometry/convex-hull.html}
\item cp-algorithms, \emph{Lattice points inside non-lattice
polygon} --- Pick's theorem's generalisation.
\newline\url{https://cp-algorithms.com/geometry/lattice-points.html}
\item cp-algorithms, \emph{Finding the nearest pair of
points} --- the strip argument written out.
\newline\url{https://cp-algorithms.com/geometry/nearest_points.html}
\item The USACO Guide, \emph{Geometry Primitives},
\url{https://usaco.guide/plat/geo-pri} --- where these
primitives get consumed at Gold/Platinum.
\item UNSW COMP4128, \emph{Computational Geometry} lecture
notes --- contest-grade slides covering this chapter plus
half-planes (file \texttt{11-comp-geom.pdf} under:)
\newline\url{https://cgi.cse.unsw.edu.au/~cs4128/25t2/lectures/}
\item M.~de Berg et al., \emph{Computational Geometry:
Algorithms and Applications} (3rd ed.), ISBN 978-3540779735
--- the textbook reference when you need the proofs.
\end{itemize}

\begin{rulebox}[Watch instead]
\begin{itemize}
\item \textbf{nptelhrd --- ``Mod-04 Lec-07 Convex Hull''},
52m35s, \url{https://www.youtube.com/watch?v=QnL0LmOO4rc} ---
covers \S16.4: hull definitions and the Graham family the
monotone chain simplifies. (Verified 2026-10-11.)
\item \textbf{Coderisland --- ``Geometric Algorithms: Graham \&
Jarvis --- Lecture 10''} (E.~Demaine, MIT 6.851),
79m37s, \url{https://www.youtube.com/watch?v=J5aJEcOr6Eo} ---
covers \S16.4 and the turn-test analysis; notation: he writes
$\texttt{ccw}$ for our \texttt{orient}. (Verified 2026-10-11.)
\item \textbf{Aaron Becker --- ``Graham Scan Tutorial: Convex
Hull of a Set of 2D Points''}, 3m06s,
\url{https://www.youtube.com/watch?v=9rQMLpQn5xQ} --- covers
\S16.4: a 3-minute animated pass over the same algorithm; watch
this first. (Verified 2026-10-11.)
\end{itemize}
\emph{Notation note:} all three use angle-sort (Graham) where we
use $(x,y)$-sort (monotone chain); the pop condition and hull
are identical.
\end{rulebox}
The orientation-first style follows cp-algorithms and
Laaksonen's geometry chapter~\cite{cpalgorithms,laaksonen2024};
hull/sweep treatments track Sannemo and the UNSW
slides~\cite{sannemo2025,usacoguide}; problem selection follows
the official CSES geometry set~\cite{cses}.
